\section{Introduction}
\label{sec:introduction}

Execution feedback enables code repair models to achieve 2.2$\times$ higher bug fix rates than global rewrites, yet most refinement methods treat feedback as a black-box signal. When a language model generates buggy code, detailed error traces---line numbers, expected versus actual values, stack traces---provide ground-truth localization that fundamentally differs from the approximate signals AI critics can offer. Understanding this difference is not merely academic: practitioners designing iterative refinement pipelines face a choice between execution sandboxes and LLM-based critics, with significant implications for infrastructure cost, latency, and repair success rates.

The surface problem is well known. Iterative code refinement with feedback improves code generation quality. Both execution-based approaches~\cite{chen2023selfdebug} and AI-based approaches~\cite{madaan2023selfrefine} have demonstrated substantial improvements over single-shot generation. However, a deeper problem emerges upon closer examination: no controlled comparison exists between these feedback types under identical experimental conditions. Prior work comparing methods such as CodeRL~\cite{le2022coderl} and Self-Refine inadvertently confounds the feedback signal with differences in model architecture, prompting strategy, and refinement loops. This leaves a critical gap: we cannot determine whether execution feedback's empirical advantage arises from ground-truth error localization or from correlated method differences.

Our key insight resolves this ambiguity. The execution feedback advantage arises not from information quantity but from information structure---error traces contain counterfactual localization (``if X were different, Y would not have failed'') that constrains the model's hypothesis space for edits. We find that 84.7\% of execution traces achieve a counterfactual information score of 0.4 or higher, with type errors (0.665) and logic errors (0.583) providing the richest diagnostic signals. This dense, actionable information enables targeted edits ($\leq$5 lines changed) that are 2.2$\times$ more likely to fix bugs than global rewrites attempting wholesale code regeneration.

Building on this insight, we present the first controlled comparison isolating feedback signals from method confounds. Our experimental design uses four conditions---execution-detailed, execution-binary, AI-critic, and random baseline---with template conversion to natural language that controls for format differences. We evaluate on CodeLlama-7B-Instruct across HumanEval and MBPP benchmarks, measuring pass@1 and refinement efficiency across 3 iterations.

Our contributions unfold as follows. First, we quantify the counterfactual information density in execution traces, providing the first systematic measurement of CF-score distributions across error types. Second, we establish the three-step causal chain---localization $\rightarrow$ targeted editing $\rightarrow$ higher fix probability---through sequential mechanism experiments. Third, we demonstrate that feedback granularity is causal: detailed execution feedback achieves 100\% refinement success on failing problems while binary pass/fail achieves only 60\%, a 40 percentage point pass@1 improvement. Finally, we report a surprising finding that challenges intuition: the execution advantage is uniform or slightly larger on simpler tasks, contrary to expectations that complex tasks would benefit more from precise localization.

We organize the remainder of this paper as follows. Section~\ref{sec:related} surveys related work on code refinement and feedback mechanisms. Section~\ref{sec:methodology} describes our methodology for isolating feedback signals. Section~\ref{sec:experiments} presents the experimental setup. Section~\ref{sec:results} reports results. Section~\ref{sec:discussion} discusses implications and limitations. Section~\ref{sec:conclusion} concludes.
